(1) 单粒子旋转系有效能为 \(\tfrac12m|\mathbf v-\Omega\times\mathbf r|^2-m\Omega^2r^2/2\)。令 \(\alpha=m\Omega^2/(2k_BT)>0\)，面积积分 \(\int_0^R2\pi r e^{\alpha r^2}dr=\pi(e^{\alpha R^2}-1)/\alpha\)，故 \(n(r)=N\alpha e^{\alpha r^2}/[\pi H(e^{\alpha R^2}-1)]\)。

(2) 局部理想气体压强 \(p(r)=n(r)k_BT\)，对 \(r\) 求导得 \(p'(r)=2\alpha r n(r)k_BT=mn(r)\Omega^2r\)。压强向外增大，与离心力指向外侧相符，并不存在体积均匀压强。

(3) 对位置配分因子取 \(\alpha\) 导数有 \(\langle r^2\rangle=\partial_\alpha\ln[(e^{\alpha R^2}-1)/\alpha]=R^2/(1-e^{-\alpha R^2})-1/\alpha\)。实验室系平均角动量 \(L_z=Nm\Omega\langle r^2\rangle\)。\(\Omega\to0\) 时 \(\langle r^2\rangle\to R^2/2\)；强转速时趋 \(R^2\)，气体被推向壁面。

(4) 平衡权重来自旋转系的 \(H-\Omega L_z\)，而非实验室能量 \(H\)。容器停止旋转后原来分子仍具有平均方位流动和非均匀密度，需要通过碰撞、流动及与壁的动量交换逐步弛豫到新平衡；不可以在突然停转瞬间直接设 \(n(r)\) 均匀。
